\documentclass[11pt,a4paper]{article}
\usepackage[margin=27mm]{geometry}
\usepackage{fontspec,amsmath,amssymb,amsthm,unicode-math}
\setmainfont{STIXGeneral}\setmathfont{STIX Two Math}
\setlength{\emergencystretch}{2em}
\newcommand{\Sym}{\operatorname{Sym}}\newcommand{\gen}[1]{\langle#1\rangle}
\newcommand{\eps}{\varepsilon}\newtheorem{proposition}{Proposition}
\title{A disconnected input to the rank and degree extension}
\author{Human Compression project}
\date{Bounded correction to the specified author edition}
\begin{document}\maketitle
\begin{abstract}
We give a string C-group on nine points satisfying the perfect-split,
fracture-graph and endpoint-movement hypotheses printed in Theorem 6.4 of
the specified Cameron--Fernandes--Leemans author edition. Its prescribed
rank and degree extension is not a C-group. The example is disconnected;
it makes no assertion against the paper's transitive stability theorem.
Every group and intersection needed for the counterexample is determined
below without relying on a finite search.
\end{abstract}

\section{The source statement and conventions}
Our source is P. J. Cameron, M. E. Fernandes and D. Leemans,
\emph{The number of string C-groups of high rank},
arXiv:2212.12723v2, 11 January 2023. All page numbers refer to that
author edition. A string group generated by involutions has ordered
involutions $\rho_0,\ldots,\rho_{r-1}$ with distant commutations.
It is a string C-group when, for every two sets of indices $I,J$,
\[
 \gen{\rho_i:i\in I}\cap\gen{\rho_j:j\in J}
       =\gen{\rho_k:k\in I\cap J}.
\]
Write $G_i=\gen{\rho_j:j\ne i}$.

The definition of a fracture graph on pp.~8--9 assumes that each $G_i$
has at least one more orbit than $G$. One chooses, for each label $i$,
an edge of $\rho_i$ joining distinct $G_i$-orbits. The resulting graph,
including isolated vertices, is a spanning forest. This definition is not
restricted to transitive groups.

For the split definition on pp.~16--17, $G_i$ has exactly one more orbit
than $G$. There is a partition $A\sqcup B$ of the point set and a unique
edge $\{a,b\}$ of $\rho_i$ crossing this partition. The split is perfect
when every earlier generator is supported on $A$ and every later generator
is supported on $B$. Write
$\rho_i=\alpha_i(a\ b)\beta_i$, with the two remaining factors supported
on their respective sides. On p.~31 the extension adds a new point $c$
and replaces this generator by the consecutive pair
\[
 \delta_i=\alpha_i(a\ c),\qquad
 \delta_{i+1}=\beta_i(c\ b),
\]
retaining earlier generators and shifting the later indices up by one.

Theorem 6.4 on pp.~32--33 says that this operation preserves the C-property
when the original marking is C, has a fracture graph, neither $a$ nor $b$
is fixed by $G_i$, and at least one of $\alpha_i,\beta_i$ is trivial.
The paragraph preceding Proposition 6.3 explicitly allows intransitive
actions in these results. We verify all these hypotheses literally.
For the old C-property we retain the standard end-parabolic criterion,
stated in source Proposition 2.1 and attributed there to
McMullen--Schulte, Proposition 2E16: a string marking with both end
parabolics C is C if their intersection is the subgroup on their common
labels.

\section{The original C-marking}
Let $\sigma=(6\ 7)$ and $\tau=(8\ 9)$, and set
\begin{equation}\label{eq:old}
 \rho_0=(1\ 2)\tau,\qquad
 \rho_1=(2\ 3)\sigma,\qquad
 \rho_2=(3\ 4)\sigma,\qquad
 \rho_3=(4\ 5).
\end{equation}
The four permutations are distinct nonidentity involutions. Every distant
pair commutes, as is seen directly from their supports. They act on
$\{1,\ldots,9\}$, with two central permutations $\sigma,\tau$ fixing
$\{1,\ldots,5\}$. In the following formulas $\eps(g)\in\mathbb F_2$
is the parity of a permutation. Groups on subsets fix the remaining points.

\begin{proposition}\label{prop:old}
The marking \eqref{eq:old} is a string C-group. Its group is
$\Sym(\{1,2,3,4,5\})\times\gen{\sigma}\times\gen{\tau}$, of order $480$.
\end{proposition}
\begin{proof}
The projection onto the first five points contains all adjacent
transpositions. Furthermore
\[
 (\rho_0\rho_1)^3=\sigma\tau,\qquad
 (\rho_2\rho_3)^3=\sigma.
\]
Thus both independent central factors belong to the group, proving the
stated full-group description.

For the intersection property put
\[
 P=\gen{\rho_0,\rho_1,\rho_2},\qquad
 Q=\gen{\rho_1,\rho_2,\rho_3},\qquad
 M=\gen{\rho_1,\rho_2}.
\]
Their actual permutation groups are
\begin{align*}
 P&=\{g\sigma^u\tau^v:
       g\in\Sym(\{1,2,3,4\}),\ \eps(g)=u+v\},\\
 Q&=\Sym(\{2,3,4,5\})\times\gen{\sigma},\\
 M&=\{g\sigma^{\eps(g)}:g\in\Sym(\{2,3,4\})\}.
\end{align*}
Here and below $u,v\in\mathbb F_2$. To prove the first formula, all three
generators satisfy its parity equation and project onto the adjacent
transpositions of $S_4$. The central element $\sigma\tau$ is already in
$P$. Hence over each of the $24$ projected permutations, $P$ contains
both possibilities permitted by that equation; these are all possibilities.
For $Q$ the projection is again $S_4$, and the second displayed cube gives
the full independent factor $\gen{\sigma}$. For $M$ the two adjacent
transpositions of $S_3$ each carry exactly one $\sigma$, so the indicated
parity graph is the group they generate.

We check that the rank-three markings of $P$ and $Q$ are C before applying
the criterion to the whole marking. The same projection and cube arguments
give
\begin{align*}
 \gen{\rho_0,\rho_1}
 &=\{g\sigma^u\tau^v:
     g\in\Sym(\{1,2,3\}),\ \eps(g)=u+v\},\\
 \gen{\rho_2,\rho_3}
 &=\Sym(\{3,4,5\})\times\gen{\sigma}.
\end{align*}
An element common to the first of these groups and $M$ has its root
permutation supported on $\{2,3\}$. Membership in $M$ sets $v=0$
and $u=\eps(g)$. Their intersection is therefore exactly
$\gen{(2\ 3)\sigma}=\gen{\rho_1}$.
Similarly, intersection of the second group with $M$ restricts the root
permutation to $\{3,4\}$ and its $\sigma$-exponent to its parity.
It is exactly $\gen{\rho_2}$.
Each group on two distinct involutions has the rank-two intersection
property. The end criterion consequently proves the C-property of both
$P$ and $Q$ with their displayed markings.

Finally an element of $P\cap Q$ has root permutation supported on
$\{2,3,4\}$ and has $v=0$ because $Q$ fixes $8$ and $9$ separately.
The equation defining $P$ then forces $u=\eps(g)$, giving precisely $M$.
A last application of the same end criterion proves the proposition.
\end{proof}

\section{Every geometric hypothesis and the failure after extension}
The full group has three orbits, namely
$\{1,2,3,4,5\}$, $\{6,7\}$ and $\{8,9\}$.
The orbits of each maximal parabolic are listed explicitly:
\[
\begin{array}{c|l}
G_0&\{1\},\ \{2,3,4,5\},\ \{6,7\},\ \{8\},\ \{9\}\\
G_1&\{1,2\},\ \{3,4,5\},\ \{6,7\},\ \{8,9\}\\
G_2&\{1,2,3\},\ \{4,5\},\ \{6,7\},\ \{8,9\}\\
G_3&\{1,2,3,4\},\ \{5\},\ \{6,7\},\ \{8,9\}.
\end{array}
\]
These are also the connected components of the respective marked permutation
graphs, so the table follows directly from \eqref{eq:old}.
Each row has at least four orbits. The four edges
$\{1,2\},\{2,3\},\{3,4\},\{4,5\}$, with labels $0,1,2,3$ respectively,
join distinct orbits in their respective rows. They form a fracture graph,
with the other four vertices isolated.

Take $i=1$, $a=2$, $b=3$ and
\[
 A=\{1,2,8,9\},\qquad B=\{3,4,5,6,7\}.
\]
The subgroup $G_1$ has exactly one more orbit than $G$.
Only the edge $\{2,3\}$ of $\rho_1$ crosses $A,B$;
the other edge $\{6,7\}$ lies within a single $G_1$-orbit.
The earlier mark $\rho_0$ is supported on $A$, and both later marks
are supported on $B$. This is a perfect split with
$\alpha_1=1$ and $\beta_1=\sigma$.
Moreover $2^{\rho_0}=1$ and $3^{\rho_2}=4$, so neither endpoint
is fixed by $G_1$.

The prescribed operation with $c=10$ produces
\[
 (\delta_0,\delta_1,\delta_2,\delta_3,\delta_4)
 =\big((1\ 2)\tau,\ (2\ 10),\ (3\ 10)\sigma,
                     \ (3\ 4)\sigma,\ (4\ 5)\big).
\]
These are again distinct involutions with the string commutations.
But the disjoint index pairs $\{1,2\}$ and $\{3,4\}$ yield
\[
 (\delta_1\delta_2)^3=\sigma
       =(\delta_3\delta_4)^3\ne1.
\]
Indeed each product is a $3$-cycle on the first five points together
with $10$, multiplied by the disjoint transposition $\sigma$.
The two generated subgroups therefore have a nontrivial intersection,
whereas the subgroup on their common index set is trivial.
The new marking is not a C-group.

\section{Scope and verification}
This counterexample disproves Theorem 6.4 in its printed intransitive
generality. It does not refute a transitive restriction or the paper's main
stability theorem, nor does it establish a corrected sufficient hypothesis.
The false general endpoint converse in source Proposition 2.4 is treated
in a separate bounded note. The source's boundary proof on p.~33 removes
an endpoint factor and regards the remaining marking as C; the present
example shows why that restricted use also needs justification.
The valid forward endpoint-extension theorem used in the admitted L010
constructions remains unaffected.

The complete accompanying GAP program verifies the literal hypotheses,
the displayed subgroup descriptions, all $256$ old and $1024$ new
intersection pairs, and the two witness words. These checks are diagnostics
of the written proof, not premises of it. No full Lean formalization,
complete L019 recovery, new conjecture answer, historical priority, or
numerical compression is claimed.
\end{document}
